\chapter{Лабораторная работа №8. Основы сетевого программирования и автоматизации.}
\label{ch:chapter 8}

\section{Цели}

Лабораторная работа помогает получить практические навыки по изучению
следующих тем:

\begin{itemize}

\item Изучение базового синтаксиса языка Python

\item Использование библиотеки telnetlib для автоматизации задач

\end{itemize}

\section{Топология сети}

У компании есть коммутатор с IP-адресом управления 192.168.56.101/24.
Необходимо написать сценарий автоматизации для просмотра текущего
конфигурационного файла устройства.

\includegraphics{lb8_1}

\section{План работы}

\begin{itemize}

\item 1. Настройка параметров Telnet: установка пароля Telnet, включение Telnet и настройка разрешения на доступ через Telnet.

\item 2. Компиляция сценария Python: вызов модуля telnetlib для входа в устройство и проверка конфигурации.

\end{itemize}

\section{Процедура конфигурирования}

\textbf{Шаг 1.} Настройка Telnet на коммутаторе.

Войдем в системный режим и зададим пароль для входа Telnet:

\begin{verbatim}
<Huawei> system-view
[Huawei] user-interface vty 0 4
[Huawei-ui-vty0-4] authentication-mode password
[Huawei-ui-vty0-4] set authentication password simple Huawei@123
[Huawei-ui-vty0-4] protocol inbound telnet
[Huawei-ui-vty0-4] user privilege level 15
[Huawei-ui-vty0-4] quit
\end{verbatim}

\textbf{Шаг 2.} Включение службы Telnet.

Включим службу Telnet, чтобы разрешить удаленный доступ:

\begin{verbatim}
[Huawei] telnet server enable
Info: The Telnet server has been enabled.
[Huawei] quit
\end{verbatim}

\textbf{Шаг 3.} Проверка доступа через Telnet с ПК.

Выполним вход с ПК в коммутатор через Telnet с помощью командной строки:

\begin{verbatim}
C:\Users\Student> telnet 192.168.56.101
Login authentication

Password: 
Info: The max number of VTY users is 5, and the number of current VTY 
users on line is 1. The current login time is 2024-01-15 21:12:57.
<Huawei>
\end{verbatim}

Как видно, функции Telnet сконфигурированы корректно.

\textbf{Шаг 4.} Написание кода на языке Python.

Создадим сценарий, который вызывает модуль telnetlib для входа в коммутатор,
выполняет команду \textbf{display current-configuration} и выводит результаты на экран:

\begin{lstlisting}[language=Python, caption=Скрипт автоматизации на Python]
import telnetlib
import time

host = '192.168.56.101'
password = 'Huawei@123'

tn = telnetlib.Telnet(host)

tn.read_until(b"Password:")
tn.write(password.encode('ascii') + b"\n")
tn.write(b'display cu \n')
time.sleep(2)

print(tn.read_very_eager().decode('ascii'))
tn.close()
\end{lstlisting}

\textbf{Шаг 5.} Запуск компилятора.

В данной лабораторной среде используется компилятор Jupyter Notebook.
Запустим написанный скрипт.

\textbf{Шаг 6.} Интерпретация кода.

Сначала импортируем необходимые модули:

\begin{verbatim}
import telnetlib   # Модуль для работы с Telnet
import time        # Модуль для работы с задержками
\end{verbatim}

Затем выполняется вход в устройство. Функция \texttt{read\_until(b"Password:")} ожидает приглашения к вводу пароля.
Функция \texttt{write()} отправляет пароль и команду \texttt{display cu}.
Метод \texttt{read\_very\_eager()} считывает вывод команды, который отображается на консоли.
В конце сеанс закрывается командой \texttt{close()}.

\textbf{Шаг 7.} Вывод командной строки.

В результате выполнения скрипта мы видим вывод конфигурации устройства:

\begin{verbatim}
Info: The max number of VTY users is 5, and the number 
 of current VTY users on line is 2. 
 The current login time is 2024-01-15 20:19:11.
<Huawei>display cu
#
sysname Huawei
#
cluster enable
ntdp enable
ndp enable
#
drop illegal-mac alarm
#
diffserv domain default
#
drop-profile default
# 
aaa
 authentication-scheme default
 authorization-scheme default
 accounting-scheme default
 domain default
 domain default_admin
 local-user admin password simple admin
 local-user admin service-type http
#
interface Vlanif1
 ip address 192.168.56.101 255.255.255.0
#
--- More ---
\end{verbatim}

\section{Вопросы}

\textbf{1. Как с помощью telnetlib настроить параметры устройства, например, IP-адрес интерфейса управления устройством?}

\textbf{Ответ:}

Для настройки параметров устройства (например, IP-адреса интерфейса) с помощью telnetlib необходимо выполнить следующие шаги:

\begin{enumerate}
    \item Установить соединение с устройством через telnetlib.
    \item Войти в систему, передав имя пользователя и пароль (или только пароль, если аутентификация настроена соответствующим образом).
    \item Перейти в системный режим командой \texttt{system-view}.
    \item Войти в режим конфигурирования интерфейса командой \texttt{interface} (например, \texttt{interface Vlanif1}).
    \item Выполнить команду настройки IP-адреса, например \texttt{ip address 192.168.1.1 24}.
    \item Сохранить конфигурацию командой \texttt{save}.
\end{enumerate}

Пример кода для настройки IP-адреса:

\begin{lstlisting}[language=Python]
import telnetlib
import time

host = '192.168.56.101'
password = 'Huawei@123'

tn = telnetlib.Telnet(host)

# Вход в систему
tn.read_until(b"Password:")
tn.write(password.encode('ascii') + b"\n")
time.sleep(1)

# Переход в системный режим и настройка IP-адреса
tn.write(b"system-view \n")
time.sleep(1)
tn.write(b"interface Vlanif1 \n")
time.sleep(1)
tn.write(b"ip address 192.168.1.1 24 \n")
time.sleep(1)
tn.write(b"quit \n")
time.sleep(1)
tn.write(b"return \n")
time.sleep(1)

# Сохранение конфигурации
tn.write(b"save \n")
time.sleep(1)
tn.write(b"y \n")
time.sleep(3)

tn.close()
\end{lstlisting}

\textbf{2. Как сохранить файл конфигурации в локальный каталог?}

\textbf{Ответ:}

Для сохранения файла конфигурации в локальный каталог можно использовать следующий подход:

\begin{enumerate}
    \item Подключиться к устройству через telnetlib и выполнить команду \texttt{display current-configuration} (или \texttt{display cu}).
    \item Считать вывод команды с помощью метода \texttt{read\_very\_eager()} или \texttt{read\_until()}.
    \item Записать полученные данные в локальный файл с помощью стандартных средств Python (функция \texttt{open()}).
\end{enumerate}

Пример кода для сохранения конфигурации:

\begin{lstlisting}[language=Python]
import telnetlib
import time

host = '192.168.56.101'
password = 'Huawei@123'

tn = telnetlib.Telnet(host)

tn.read_until(b"Password:")
tn.write(password.encode('ascii') + b"\n")
time.sleep(1)

tn.write(b"display current-configuration \n")
time.sleep(3)

config = tn.read_very_eager().decode('ascii')
tn.close()

# Сохранение конфигурации в локальный файл
with open('router_config.txt', 'w') as f:
    f.write(config)
    
print("Конфигурация сохранена в файл router_config.txt")
\end{lstlisting}